\section{\slowlab}
\label{sec:regime}
\label{sec:env}

\subsection{Regime and problem formulation}
\label{sec:conditions}
\label{sec:formulation}

\slowlab{} targets sequential experimental design under four conditions that hold jointly.
\textbf{Slow}: what can be learnt per unit of calendar time is bounded by feedback latency
\emph{and} by a hard ceiling on how many trials run concurrently---a facility with a
thousand chambers and a four-month cycle is not slow in the sense that binds. That ceiling
is \emph{tiered}, because some factors are set for a whole block of units and others per
unit, so replication and coverage compete for the same physical resource.
\textbf{Noisy}: a single unit does not settle a question, which is what turns replication,
blocking and allocation into decisions rather than formalities. \textbf{Costly}: a trial
consumes resources measured in the objective's own units, so exploration is priced rather
than free. \textbf{Irreversible}: a submitted design cannot be rolled back---no branch, no
checkpoint, no revert---so the only remedy for a bad design is a new one in a later cycle,
and the units it occupied are gone. Table~\ref{tab:landscape} places existing evaluation
environments against these conditions. We instantiate the regime in controlled-environment
agriculture (\S\ref{sec:domain}).

\begin{table}[t]
\centering\footnotesize
\setlength{\tabcolsep}{3.5pt}
\begin{tabular}{@{}llccc@{}}
\toprule
& & \textbf{Agent / design} & \textbf{Chemistry \& materials} & \textbf{\slowlab} \\
& & {\scriptsize BoxingGym} & {\scriptsize A-Lab, Coscientist} & {\scriptsize (this work)} \\
\midrule
\multirow{2}{*}{\textbf{Slow}}
  & feedback latency        & seconds--minutes & hours--days & \textbf{210 days} \\
  & parallel width          & unbounded & high (96/384-well) & \textbf{hard-capped, tiered} \\
\addlinespace[2pt]
\textbf{Noisy}
  & signal-to-noise         & deterministic & high & \textbf{low (plant-to-plant)} \\
\addlinespace[2pt]
\multirow{2}{*}{\textbf{Costly}}
  & cost of one trial       & $\approx 0$ & reagents, instrument time & \textbf{a crop, priced in the objective} \\
  & cost of a retry         & $\approx 0$ & low--moderate & \textbf{one growing season} \\
\addlinespace[2pt]
\textbf{Irreversible}
  & can a design be undone? & fully (reset) & largely (re-run) & \textbf{no} \\
\bottomrule
\end{tabular}
\caption{Evaluation environments against the four conditions of \S\ref{sec:regime}. Each row
is one measurable property of the condition naming it. The closest prior benchmark differs
on every row.}
\label{tab:landscape}
\end{table}

\paragraph{Problem formulation.} An instance is a parameter vector $\theta \sim P_\Theta$,
unknown to the agent, determining a response $f_\theta : \mathcal{X} \to \mathbb{R}$ on a
factor space $\mathcal{X} \subseteq \mathbb{R}^{d}$ and measured in the objective's own
units. Over $R$ rounds the agent commits a \emph{design} $D_r \in \mathcal{D}$ --- a finite
set of treatments in $\mathcal{X}$ together with their assignment to the physical trial
slots $\mathcal{U}_r$ --- and states a recommendation $\hat{x}_r \in \mathcal{X}$. The
admissible set $\mathcal{D}$ belongs to the facility rather than to the task;
\S\ref{sec:facility} gives it for our instantiation. Its one structural role here is that it
is not all finite multisets: an agent cannot buy precision and coverage with the same slots.

Committing holds $\mathcal{U}_r$ for a cycle of length $\Delta$, after which each slot
returns one observation
\begin{equation}
  y_u \;=\; f_\theta\!\left(x_u;\, \phi_u\right) \;+\; \varepsilon_u,
  \qquad \varepsilon \sim \mathcal{N}\!\big(0, \Sigma(D_r)\big),
  \label{eq:obs}
\end{equation}
where $x_u$ is the treatment assigned to slot $u$ and $\phi_u$ an independent per-slot draw
of the response model's own parameters. Two features of \eqref{eq:obs} carry \textbf{noisy}.
Drawing $\phi$ per slot rather than adding a scalar error makes the variance heteroscedastic
in $x$ and structured by the mechanism; and $\Sigma(D_r)$ is not diagonal, so \emph{where} a
treatment is placed changes what the round can resolve, not only which treatments were
chosen. Appendix~\ref{app:eig} gives $\Sigma$ for our instantiation.

The other three conditions restrict the action. Writing $\mathcal{F}_r = \sigma\big(\{\,y_u
: u \in \mathcal{U}_1 \cup \cdots \cup \mathcal{U}_{r-1}\,\}\big)$ for everything returned
before round $r$, and $\mathcal{A}_r$ for the entire action set available at round $r$,
\begin{equation}
  D_r,\ \hat{x}_r \ \text{ are } \mathcal{F}_r\text{-measurable},
  \qquad\qquad
  \mathcal{A}_r \;=\; \big\{(D_r,\ \hat{x}_r)\big\}.
  \label{eq:action}
\end{equation}
The first is \textbf{slow} stated as measurability rather than as sensor deprivation:
information about a committed slot may well exist before its cycle ends, but no action is
available on it, and Appendix~\ref{app:defects} names the two actions that would separate
the two readings. The second is \textbf{irreversible}: the action set contains commit and
recommend and nothing else, from which it follows that no operation shrinks $\mathcal{F}_r$,
releases a slot early, or refunds one already committed. \textbf{Costly} is carried not by a
constraint but by the objective, which prices a trial in the units it is trying to maximise,
so that a campaign has a cost beside the quality of its answer (\S\ref{sec:eval}).

\paragraph{The problem.} Given $\mathcal{F}_r$ and nothing else, the agent chooses
$D_1, \dots, D_R$ and finally commits to $\hat{x}_R$. It is judged on how far that
recommendation falls short of the instance's own optimum,
\begin{equation}
  \mathcal{R} \;=\; f_\theta(x^\star_\theta) - f_\theta(\hat{x}_R),
  \qquad x^\star_\theta \in \arg\max_{x \in \mathcal{X}} f_\theta(x),
  \label{eq:problem}
\end{equation}
and, separately, on what the campaign spent reaching it: every slot it committed is charged
the same shortfall, at the treatment that slot was given, for a full cycle
\eqref{eq:cumregret}. Neither quantity
is observable to the agent --- $\theta$, $f_\theta$ and $x^\star_\theta$ are all withheld ---
and neither can be estimated from $\mathcal{F}_r$ without the design that generated it.

The two pull against each other, and the constraints are what stop the usual resolution.
Reducing \eqref{eq:problem} calls for slots spent on treatments the agent expects to be
poor, since those are the ones that discriminate; that is exactly what the second charge
penalises. In an environment where a trial is a function call, the tension is nominal: try
something, look, retract it. Here \eqref{eq:action} removes retraction, \eqref{eq:obs}
means one slot does not settle whether a treatment was poor, and $\mathcal{D}$ means the
slots spent discriminating are slots not spent confirming. A campaign therefore occupies
the facility for $R\Delta$ and consists of $R$ ordered, unrepeatable commitments, each made
before the information that would have justified it exists. \S\ref{sec:eval} states both
quantities precisely and what they are read against; Appendix~\ref{app:worked} works one
round through in full.

\subsection{The environment}
\label{sec:environment}

\slowlab{} is a harness: the components below are defined over a domain rather than within
one, and \S\ref{sec:domain} supplies the first domain that fills them in;
Figure~\ref{fig:overview} on page~\pageref{fig:overview} is the loop this
subsection specifies.


\subsubsection{Interaction protocol}
\label{sec:episode}

Figure~\ref{fig:protocol} is the whole interface. \texttt{validate\_design} is free
and unlimited; \texttt{submit\_design} commits the named slots for one cycle;
\texttt{advance} runs the cycle and returns one value per slot; the interim
recommendation records what the agent would advise if the campaign stopped there,
which supplies the regret trajectory. Budget is accounted in slot-cycles. The
recommendation is elicited, never consumed: it costs no budget, returns nothing, and
enters no later design --- an agent that revises it every round is not penalised for
having been wrong earlier. Design and recommendation are separate submissions, which is
what allows an agent to run a good campaign and then answer badly, and \S\ref{sec:eig} to
score the two apart.

\begin{figure}[t]
\centering
\footnotesize
\setlength{\fboxsep}{5pt}
%
\begin{minipage}[t]{0.225\linewidth}
\fbox{\begin{minipage}[t][45mm][t]{\dimexpr\linewidth-2\fboxsep-2\fboxrule}
\centering\textbf{\slowlab{} API}\par\vspace{3pt}
\raggedright\ttfamily\scriptsize
\pykw{class} \pycls{SlowLabEnv}:\\
\hspace*{1em}\pyfn{validate\_design}\\
\hspace*{1em}\pyfn{submit\_design}\\
\hspace*{1em}\pyfn{advance}\\
\hspace*{1em}\pyfn{observations}\\
\hspace*{1em}\pyfn{submit\_recommendation}\\[4pt]
\pykw{class} \pycls{Design}:\\
\hspace*{1em}treatments\\
\hspace*{1em}allocation\\
\hspace*{1em}randomization\_seed\\[4pt]
\pykw{class} \pycls{Task}:\\
\hspace*{1em}factors, lev(\,$\cdot$\,)\\
\hspace*{1em}R, |U\_r|, $\Delta$
\end{minipage}}
\end{minipage}
\hfill
\begin{minipage}[t]{0.335\linewidth}
\fcolorbox{black}{blue!4}{\begin{minipage}[t][45mm][t]{\dimexpr\linewidth-2\fboxsep-2\fboxrule}
\centering\textbf{One campaign}\par\vspace{3pt}
\raggedright\ttfamily\scriptsize
\tikzmark{a}env = \pycls{SlowLabEnv}(task, seed)\\[2pt]
\pykw{for} r \pykw{in} \pykw{range}(R):\\
\hspace*{1em}\tikzmark{b}D = agent.\pyfn{design}(env)\\
\hspace*{1em}\tikzmark{c}env.\pyfn{validate\_design}(D)\\
\hspace*{1em}\tikzmark{d}env.\pyfn{submit\_design}(D)\\
\hspace*{1em}\tikzmark{e}env.\pyfn{observations}() $\rightarrow$ []\\
\hspace*{1em}\tikzmark{f}env.\pyfn{advance}()\\
\hspace*{1em}\tikzmark{g}x\_r = agent.\pyfn{recommend}(env)\\[2pt]
\tikzmark{h}env.\pyfn{submit\_recommendation}(x\_R)\tikzmark{edge}
\end{minipage}}
\end{minipage}
\hfill
\begin{minipage}[t]{0.385\linewidth}
\fbox{\begin{minipage}[t][45mm][t]{\dimexpr\linewidth-2\fboxsep-2\fboxrule}
\centering\textbf{What the environment does}\par\vspace{3pt}
\raggedright\scriptsize
\tikzmark{A}the seed draws a \emph{site} $\theta\sim P_\Theta$, not a noise realisation\\[2.5pt]
\tikzmark{C}free and unlimited; feasibility only, and a rejection costs no budget\\[2.5pt]
\tikzmark{D}\textbf{the irreversible step}: $|\mathcal{U}_r|$ slots held for a full cycle\\[2.5pt]
\tikzmark{E}nothing is visible until the cycle ends\\[2.5pt]
\tikzmark{F}a cycle of length $\Delta$ passes; one value per committed slot\\[2.5pt]
\tikzmark{G}the interim recommendation, giving the regret trajectory\\[2.5pt]
\tikzmark{H}ends the episode; returns $\bar{\mathcal{R}}$ and $\mathcal{R}_{\mathrm{cum}}$
\end{minipage}}
\end{minipage}
%
\begin{tikzpicture}[overlay, remember picture,
   >={Stealth[length=3.2pt]}, thin, gray!70]
  \foreach \src/\dst in {a/A, c/C, d/D, e/E, f/F, g/G, h/H}
    \draw[->] let \p1=(pic cs:\src), \p2=(pic cs:edge), \p3=(pic cs:\dst) in
      (\x2+3pt, \y1+1.2pt) -- (\x3-2.5pt, \y3+1.2pt);
\end{tikzpicture}
\caption{\textbf{The interaction protocol.} \emph{(left)} What an agent touches: an
environment exposing five calls, a \texttt{Design} naming treatments and assigning them to
named slots, and a \texttt{Task} fixing the factors, their control levels, the number of
rounds, the width cap and the cycle length. \emph{(centre)} One campaign. \emph{(right)}
What each call costs. Checking is free and a rejection carries a code and a detail, so an
agent may revise without penalty; committing is the step that cannot be undone. Between
\texttt{submit\_design} and \texttt{advance} the observation list stays empty --- a
regression test guards it --- which is \eqref{eq:action} in code.}
\label{fig:protocol}
\end{figure}

\subsubsection{Action space}
\label{sec:action}

A submission is a \texttt{Design}: named treatments, each a full vector of factor values; an
allocation mapping treatments to specific units; and a declared randomisation seed. The
agent chooses \emph{what} to test, \emph{how many units} each treatment receives, and
\emph{which} units those are---replication and allocation are decisions rather than
defaults. It does not choose what to measure, when to measure, or whether to intervene
mid-cycle; a submitted design runs to completion.

\subsubsection{Facility and control levels}
\label{sec:facility}

A facility is $C$ blocks of $L$ units. Every factor carries a \emph{control level}: a
\textsc{block}-level factor takes one value for a whole block, a \textsc{unit}-level factor
is set per unit. The level is a physical property of the apparatus, not a modelling choice,
and a design assigning two different block-level values inside one block describes an
experiment that cannot be performed.

Two consequences follow, and we learned both by implementation. With $k$ block-level
factors, the number of distinct block-level \emph{combinations}---not levels per
factor---bounds the number of whole-plot treatments by $C$. And if a combination occupies
exactly one block, its effect is confounded with that block's, and unit-level replicates
cannot be spread. A valid design is therefore a \textbf{split-plot}: block-level factors as
whole plots replicated across at least two blocks, unit-level factors as subplots. This
structure is what makes allocation a decision, and it is absent from any environment in
which an experiment is a function call.

\subsubsection{Forward model and instances}
\label{sec:world}

The environment takes a published forward model $f_\theta$ mapping a factor vector to a mean
response, and a parameter distribution $P_\Theta$ describing how that model's coefficients
vary across sites in the literature. A seed selects an \emph{instance} by drawing
$\theta \sim P_\Theta$: a different problem, not the same problem observed again. The
response surface and the location of its optimum are therefore fixed by the published model
rather than by us.

Within an instance, each unit receives its own draw of the model's parameters rather than an
additive noise term. Unit-to-unit variation \emph{is} parameter variation, and modelling it
this way yields heteroscedastic, temporally correlated variability without further
assumptions. The nuisance effects of \S\ref{sec:conditions} are what an allocation must avoid
confounding with.

\subsubsection{Objective and cost}
\label{sec:costmodel}

The objective is a rate: revenue per unit area per unit time, less the full cost of the
resources the treatment consumed, $f_\theta(x) = \big(\text{revenue}(x) -
\text{cost}(x)\big) / \Delta$. Pricing is not decoration. Under a pure-yield objective a
forward model is typically monotone in its resource inputs, so their optima sit on the
boundary and those dimensions are trivial; charging for them is what produces interior
optima. Expressing the cost in the objective's own units is also what makes the
\emph{Costly} measure interpretable: a campaign's expenditure is directly comparable to
what it was trying to maximise.

The costs enter at weight one, and this is enforced rather than chosen. An earlier version
carried a multiplier $\lambda$ on the cost term, declared as a task-design parameter. It
was not one: it set where every optimum sat, and at $\lambda = 0.25$ it pushed four of six
factors onto their rails. The environment now raises if $\lambda \neq 1$
(Appendix~\ref{app:defects}). What \textsc{Transfer} moves mid-episode is a price in the
cost model, not a weight on it, so the objective stays the profit a grower would actually
bank in both halves of that task.

\subsubsection{Feasibility checking}
\label{sec:checking}

\texttt{validate\_design} checks \emph{mechanical} feasibility only---control levels, unit
conflicts, budget, factor ranges---and is free and unlimited; rejections cost no budget but
are counted. The design measure of \S\ref{sec:eig} is computed at scoring time and
cannot be called by the agent. Separating the two has a purpose: because feasibility is
free to check, a design that passes it yet is scientifically poor is unambiguously a
reasoning error rather than an interface error. Anything requiring a statistical
assumption---assumed variances, power targets---belongs to the tool layer, an experimental
condition to be ablated (Appendix~\ref{app:refs}), not a free service of the environment.

\subsubsection{Language-model interface}
\label{sec:harness}

A language model reaches the same interface through a harness that serialises the round. It
receives the facility layout and unit naming, the budget and rounds remaining, each factor's
range and control level, the hierarchy constraint, every observation so far with the settings
that produced it, the currently available units, and a JSON schema; it returns treatments, an
allocation, a randomisation seed, and the recommendation it would give if the campaign stopped
now, which supplies the trajectory in \eqref{eq:regret}.

The prompt states constraints and never gives advice: the words replication, randomisation,
blocking, split-plot, pure error and confounding do not appear, since those name the
capability under test, and a regression test enforces their absence. Failures are counted
separately---a reply that will not parse from a design that is mechanically infeasible---both
fed back verbatim and retried up to three times, and neither consuming budget, consistent with
\texttt{validate\_design} being free. ``Cannot emit valid JSON'' and ``cannot design an
experiment'' are therefore different numbers. Exhausting the retries forfeits that round and
the campaign continues, so one unparseable reply is not as costly as a failed campaign.

\subsubsection{Configuration and extension}
\label{sec:interfaces}

Configurations are declarative: factors and their control levels, facility shape, rounds,
per-round width, variance components, cost weight. Adding one requires no environment code.
Agents implement a single \texttt{run(env)} method and are registered by name; the reference
strategies of Appendix~\ref{app:refs} and the harness above use the same interface, and the
harness needs only a callable mapping a message list to a string, so no vendor SDK is
required. A new domain supplies $f_\theta$, $P_\Theta$, a cost function in the
objective's units, and a facility whose units are not interchangeable
(Appendix~\ref{app:admit}).
